% --- Language & Font Management (XeLaTeX) ---
\usepackage{fontspec}
\usepackage{polyglossia}

% Set default language to Russian with modern hyphenation
\setdefaultlanguage{russian}
\setotherlanguage{english}

% Font Configuration: Using standard Windows fonts for compatibility
% Ligatures=TeX enables standard TeX ligatures (-- for en-dash, etc.)
\setmainfont{Times New Roman}[
    Ligatures=TeX,
    Scale=1.0
]
\setsansfont{Arial}[
    Ligatures=TeX,
    Scale=MatchLowercase
]
\setmonofont{Consolas}[
    Scale=MatchLowercase,
    Contextuals=Alternate % Better coding ligatures if supported
]
\newfontfamily\codefont{Consolas}[Scale=MatchLowercase]

% Microtypography: Improves justification and readability significantly
\usepackage{microtype}
\emergencystretch=3em % Fixes text overflow by allowing more spacing flexibility
\sloppy % Relaxes line breaking rules to avoid overflows

% --- Paragraphs & Spacing (GOST) ---
\usepackage{indentfirst} % Indent first paragraph of sections
\usepackage{setspace}
\onehalfspacing % 1.5 line spacing
\setlength{\parindent}{1.25cm} % Standard GOST indent
\setlength{\parskip}{0pt} % No extra space between paragraphs

% --- Mathematics ---
\usepackage{amsmath}
\usepackage{amssymb}
% \usepackage{unicode-math} % Uncomment if you want to use OpenType math fonts

% --- Colors & Graphics ---
\usepackage[table]{xcolor}
\usepackage{graphicx}
\graphicspath{{images/}{chapters/python/images/}}
\usepackage{float}
\usepackage{tikz}
\usepackage{tikz-3dplot}
\usetikzlibrary{arrows.meta, quotes, 3d, calc, shapes.geometric, shadows, positioning, decorations.markings, fadings, patterns, shadings, shapes}

% Custom fadings for TikZ
\tikzfading[name=fade out, inner color=transparent!0, outer color=transparent!100]

% Define Modern Color Palette (Geoskan Brand)
\definecolor{GeoskanBlue}{RGB}{0, 91, 187}       % Primary Brand Color
\definecolor{GeoskanDarkBlue}{RGB}{22, 43, 58}    % Dark Text / Headers
\definecolor{GeoskanOrange}{RGB}{255, 140, 0}     % Accent / Highlight
\definecolor{GeoskanLightGray}{RGB}{245, 245, 245}% Backgrounds
\definecolor{CodeGreen}{RGB}{0, 150, 0}           % Comments
\definecolor{CodePurple}{RGB}{128, 0, 128}        % Keywords
\definecolor{CodeBack}{RGB}{250, 250, 250}        % Code Background

% Additional Colors for Drone Illustration
\definecolor{GeoskanBlack}{RGB}{30, 30, 35}
\definecolor{GeoskanDarkGrey}{RGB}{45, 45, 50}
\definecolor{GeoskanGrey}{RGB}{100, 100, 105}
\definecolor{GeoskanLightGrey}{RGB}{180, 180, 185}
\definecolor{SkyTop}{RGB}{100, 180, 255}
\definecolor{SkyBottom}{RGB}{200, 230, 255}
\definecolor{GroundGreen}{RGB}{34, 139, 34}

% --- Bibliography (BibLaTeX) ---
\usepackage[
    backend=biber,
    style=gost-numeric,
    language=auto,
    autolang=other,
    sorting=none
]{biblatex}
\addbibresource{styles/references.bib}

% --- Page Layout (GOST 7.32-2017) ---
\usepackage{geometry}
\geometry{
    a4paper,
    top=20mm,
    bottom=20mm,
    left=30mm,
    right=15mm,
    headheight=15mm,
    footskip=15mm
}

% --- Headers & Footers (scrlayer-scrpage) ---
\usepackage[automark, headsepline]{scrlayer-scrpage}

% Clear default styles
\clearpairofpagestyles

% Header: Left = Section Name, Right = Book Title
\ihead{\headmark}
\ohead{\textbf{\textcolor{GeoskanBlue}{ГЕОСКАН ПИОНЕР}}}

% Footer: Center = Page Number (GOST usually centers page numbers)
\cfoot{\pagemark}

% Font styles for headers/footers
\setkomafont{pageheadfoot}{\sffamily\small\color{GeoskanDarkBlue}}
\setkomafont{pagenumber}{\sffamily\bfseries\color{GeoskanBlue}}

% --- Typography & Section Styling (KOMA-Script) ---
% Section (Level 1) - Left aligned, bold, 14pt+
\setkomafont{section}{\sffamily\Large\bfseries\color{GeoskanBlue}}
% Subsection (Level 2)
\setkomafont{subsection}{\sffamily\large\bfseries\color{GeoskanDarkBlue}}
% Subsubsection (Level 3)
\setkomafont{subsubsection}{\sffamily\normalsize\bfseries\color{GeoskanOrange}}
% Paragraph (Level 4)
\setkomafont{paragraph}{\sffamily\bfseries\color{GeoskanDarkBlue}}

% Customizing Section Numbering format if needed
% \renewcommand*{\sectionformat}{\color{GeoskanBlue}\thesection\autodot\enskip}

% --- Lists & Tables ---
\usepackage{enumitem}
\setlist{nosep} % Compact lists by default
% GOST lists: dashes or numbers.
\setlist[itemize]{label=\textcolor{GeoskanBlue}{\textbullet}, leftmargin=*, itemsep=0pt, parsep=0pt}
\setlist[enumerate]{label=\textcolor{GeoskanBlue}{\arabic*.}, leftmargin=*, font=\bfseries, itemsep=0pt, parsep=0pt}
\setlist[description]{font=\bfseries\color{GeoskanDarkBlue}, itemsep=0pt, parsep=0pt}

\usepackage{booktabs} % Professional table rules
\usepackage{tabularx} % Auto-width tables
\usepackage{multirow}
\usepackage{caption} 
\captionsetup{font=small, labelfont={bf,color=GeoskanBlue}} % GOST caption style
\usepackage{longtable} % For long tables in appendix

% --- Code Listings (Minted) ---
\usepackage[outputdir=../build]{minted}
\usemintedstyle{friendly}
\setminted{
    fontsize=\small,
    breaklines=true,
    frame=none,
    baselinestretch=1.0,
    linenos=false
}

% --- Boxes & Environments (tcolorbox) ---
\usepackage{tcolorbox}
\usepackage{etoolbox} % For \ifstrempty
\tcbuselibrary{skins, breakable, minted}
\usepackage{fontawesome5} % Icons
\usepackage{textcomp} % For upquote
\usepackage{accsupp} % Accessibility support
\usepackage{attachfile2} % File attachments

% --- Copy-paste friendly code listings -------------------------------------
%
% Problem: the PDF text layer produced by minted/fancyvrb does not reproduce
% the original source file when the reader copies a listing out of the PDF:
%   * leading indentation is lost (spaces are TeX glue, not glyphs),
%   * the hyphen glyph is mapped to U+2010 instead of ASCII U+002D,
%   * a long line that visually wraps is copied as two lines and the
%     fvextra continuation symbol ",->" leaks into the copied text.
% All of that breaks the code as soon as it is pasted into the project's
% Monaco based web editor (and into any Lua/Python tool chain).
%
% Fix: read the real source file, and attach its exact text, line by line,
% to the typeset lines as a PDF /Span with /ActualText (accsupp).  Text
% extractors must honour /ActualText, so the copied text becomes the literal
% file content.  Nothing about the visual output changes.
%
% The ActualText deliberately carries NO line terminator: PDF viewers and
% extractors insert their own newline between two typeset lines, so adding
% U+000A here would produce doubled blank lines on paste.  Source lines that
% contain no visible glyph (empty or whitespace-only) get no span at all,
% because a marked-content span without glyphs confuses extractors.
%
% The previous \FV@Space patch (a single space wrapped in ActualText) was a
% no-op: with breaklines=true fvextra redefines \FV@Space when the listing
% environment starts, so the patch never reached the output.  It is replaced
% by the per-line mechanism below, which covers spaces as part of the line.

% accsupp's dvipdfm driver (the one XeTeX selects) emits BDC/EMC through the
% "pdf:content" special, which dvipdfmx wraps in "q <translate to the current
% point> cm ... Q".  Text extractors that compute the span's position at EMC
% time therefore see the span shifted to the right by the width of the line,
% which makes poppler place the whole listing in a bogus column and silently
% drop every listing after the first one on a page.  "pdf:code" injects the
% same operators with no transformation, which is what BDC/EMC need.
\makeatletter
\def\ACCSUPP@bdc{\special{pdf:code \ACCSUPP@span BDC}}
\def\ACCSUPP@emc{\special{pdf:code EMC}}
\makeatother

\ExplSyntaxOn
\seq_new:N  \g__codeacc_hex_seq
\int_new:N  \l__codeacc_index_int
\int_new:N  \l__codeacc_pending_int
\ior_new:N  \g__codeacc_ior
\bool_new:N \g__codeacc_active_bool
\tl_new:N   \l__codeacc_line_tl
\str_new:N  \l__codeacc_hex_str

% Store one raw source line as UTF-16BE hex (or an empty entry for a line
% that will not produce any glyph on the page).
\cs_new_protected:Npn \__codeacc_store_line:n #1
  {
    \tl_set:Nx \l__codeacc_line_tl { \tl_trim_spaces:n {#1} }
    \tl_if_empty:NTF \l__codeacc_line_tl
      {
        \int_incr:N \l__codeacc_pending_int
        \seq_gput_right:Nn \g__codeacc_hex_seq { }
      }
      {
        \str_set_convert:Nnnn \l__codeacc_hex_str {#1} { } { utf16be/hex }
        \seq_gput_right:Nx \g__codeacc_hex_seq { \str_use:N \l__codeacc_hex_str }
        \int_zero:N \l__codeacc_pending_int
      }
  }

% Read the whole source file once, before \inputminted typesets it.
\cs_new_protected:Npn \__codeacc_load:n #1
  {
    \seq_gclear:N \g__codeacc_hex_seq
    \bool_gset_false:N \g__codeacc_active_bool
    \file_if_exist:nT {#1}
      {
        \int_zero:N \l__codeacc_pending_int
        \ior_open:Nn \g__codeacc_ior {#1}
        \ior_str_map_inline:Nn \g__codeacc_ior { \__codeacc_store_line:n {##1} }
        \ior_close:N \g__codeacc_ior
        \bool_gset_true:N \g__codeacc_active_bool
      }
  }

% fvextra calls \FancyVerbFormatLine more than once per line (once inside a
% measuring \sbox), so the source line is looked up through fancyvrb's own
% FancyVerbLine counter instead of a counter of our own.
\cs_new_protected:Npn \__codeacc_wrap_line:n #1
  {
    \bool_if:NTF \g__codeacc_active_bool
      {
        \int_set:Nn \l__codeacc_index_int { \value{FancyVerbLine} }
        \str_set:Nx \l__codeacc_hex_str
          { \seq_item:Nn \g__codeacc_hex_seq { \l__codeacc_index_int } }
        \str_if_empty:NTF \l__codeacc_hex_str
          { #1 }
          {
            \exp_args:Nx \BeginAccSupp
              { method=hex,unicode,ActualText= \str_use:N \l__codeacc_hex_str }
            #1
            \EndAccSupp { }
          }
      }
      { #1 }
  }

% \CodeAccInputMinted{<language>}{<path relative to tex/>}
\NewDocumentCommand \CodeAccInputMinted { m m }
  {
    \__codeacc_load:n {#2}
    \group_begin:
      \cs_set_eq:NN \FancyVerbFormatLine \__codeacc_wrap_line:n
      \inputminted [ fontsize = \small ] {#1} {#2}
    \group_end:
  }
\ExplSyntaxOff

% 1. Code File Environment (Modern Flat Design)
\NewDocumentEnvironment{codefile}{O{python}m}{
    \begin{tcolorbox}[
        enhanced,
        colback=CodeBack,
        colframe=GeoskanBlue,
        coltitle=GeoskanDarkBlue,
        boxrule=0.5pt,
        leftrule=4pt, % Thicker accent line
        arc=2mm,
        outer arc=2mm,
        left=6mm, right=3mm, top=3mm, bottom=3mm,
        breakable,
        title={\sffamily\bfseries\small \faFileCode\ \detokenize{#2}},
        fonttitle=\sffamily\bfseries\small,
        attach boxed title to top left={xshift=4mm, yshift=-3mm},
        boxed title style={colback=white, colframe=white, boxrule=0pt},
        shadow={2mm}{-2mm}{0mm}{black!10}
    ]
    \CodeAccInputMinted{#1}{../code/#2}
    \end{tcolorbox}
}{}

% 1a. Code File Environment with Custom Title
\NewDocumentEnvironment{codefiletitle}{O{python}mm}{
    \begin{tcolorbox}[
        enhanced,
        colback=CodeBack,
        colframe=GeoskanBlue,
        coltitle=GeoskanDarkBlue,
        boxrule=0.5pt,
        leftrule=4pt, % Thicker accent line
        arc=2mm,
        outer arc=2mm,
        left=6mm, right=3mm, top=3mm, bottom=3mm,
        breakable,
        title={\sffamily\bfseries\small \faFileCode\ \detokenize{#2}},
        fonttitle=\sffamily\bfseries\small,
        attach boxed title to top left={xshift=4mm, yshift=-3mm},
        boxed title style={colback=white, colframe=white, boxrule=0pt},
        shadow={2mm}{-2mm}{0mm}{black!10}
    ]
    \CodeAccInputMinted{#1}{../code/#3}
    \end{tcolorbox}
}{}

% 2. Info Box (Blue)
\newtcolorbox{infobox}[1][Информация]{
    enhanced,
    colback=blue!3!white,
    colframe=GeoskanBlue,
    coltitle=GeoskanBlue,
    boxrule=0.5pt,
    leftrule=4pt,
    arc=2mm,
    title={\sffamily\bfseries \faInfoCircle\ #1},
    fonttitle=\sffamily\bfseries,
    attach boxed title to top left={xshift=4mm, yshift=-3mm},
    boxed title style={colback=white, colframe=white, boxrule=0pt},
    breakable,
    top=4mm,
    shadow={2mm}{-2mm}{0mm}{black!10}
}

% 3. Warning Box (Red)
\newtcolorbox{warnbox}[1][Внимание]{
    enhanced,
    colback=red!3!white,
    colframe=red!70!black,
    coltitle=red!70!black,
    boxrule=0.5pt,
    leftrule=4pt,
    arc=2mm,
    title={\sffamily\bfseries \faExclamationTriangle\ #1},
    fonttitle=\sffamily\bfseries,
    attach boxed title to top left={xshift=4mm, yshift=-3mm},
    boxed title style={colback=white, colframe=white, boxrule=0pt},
    breakable,
    top=4mm,
    shadow={2mm}{-2mm}{0mm}{black!10}
}

% 4. Task Box (NO FRAMES - Just numbered text)
\newcounter{taskcounter}
\setcounter{taskcounter}{0}

\newenvironment{taskbox}[1][]{%
    \par\bigskip% Vertical space before
    \refstepcounter{taskcounter}% Increment
    \noindent\textbf{\large\sffamily\color{GeoskanBlue} Задание \thetaskcounter\ifstrempty{#1}{}{: #1}}% Title
    \par\medskip\noindent\ignorespaces% Space after title and start text
}{%
    \par\bigskip% Vertical space after
}

% 5. Hint Box (Orange)
\newtcolorbox{hintbox}[1][Подсказка]{
    enhanced,
    colback=orange!5!white,
    colframe=GeoskanOrange,
    coltitle=GeoskanOrange,
    boxrule=0.5pt,
    leftrule=4pt,
    arc=2mm,
    title={\sffamily\bfseries \faLightbulb\ #1},
    fonttitle=\sffamily\bfseries,
    attach boxed title to top left={xshift=4mm, yshift=-3mm},
    boxed title style={colback=white, colframe=white, boxrule=0pt},
    breakable,
    top=4mm,
    shadow={2mm}{-2mm}{0mm}{black!10}
}

% 6. Rule Box (Dark Blue - Rules/Conventions)
\newtcolorbox{rulebox}[1][Правило]{
    enhanced,
    colback=GeoskanDarkBlue!5!white,
    colframe=GeoskanDarkBlue,
    coltitle=GeoskanDarkBlue,
    boxrule=0.5pt,
    leftrule=4pt,
    arc=2mm,
    title={\sffamily\bfseries \faGavel\ #1},
    fonttitle=\sffamily\bfseries,
    attach boxed title to top left={xshift=4mm, yshift=-3mm},
    boxed title style={colback=white, colframe=white, boxrule=0pt},
    breakable,
    top=4mm,
    shadow={2mm}{-2mm}{0mm}{black!10}
}

% 7. Alert Box (Red/Important)
\newtcolorbox{alertbox}[1][Важно]{
    enhanced,
    colback=red!5!white,
    colframe=red!80!black,
    coltitle=red!80!black,
    boxrule=0.5pt,
    leftrule=4pt,
    arc=2mm,
    title={\sffamily\bfseries \faExclamationCircle\ #1},
    fonttitle=\sffamily\bfseries,
    attach boxed title to top left={xshift=4mm, yshift=-3mm},
    boxed title style={colback=white, colframe=white, boxrule=0pt},
    breakable,
    top=4mm,
    shadow={2mm}{-2mm}{0mm}{black!10}
}

% 8. Question Box (Green/Questions)
\newtcolorbox{questionbox}[1][Вопрос]{
    enhanced,
    colback=green!5!white,
    colframe=green!60!black,
    coltitle=green!60!black,
    boxrule=0.5pt,
    leftrule=4pt,
    arc=2mm,
    title={\sffamily\bfseries \faQuestionCircle\ #1},
    fonttitle=\sffamily\bfseries,
    attach boxed title to top left={xshift=4mm, yshift=-3mm},
    boxed title style={colback=white, colframe=white, boxrule=0pt},
    breakable,
    top=4mm,
    shadow={2mm}{-2mm}{0mm}{black!10}
}

% 9. Simple Code Box
\newtcolorbox{codebox}[1][]{
    enhanced,
    colback=CodeBack,
    colframe=gray!30,
    boxrule=0.5pt,
    arc=2mm,
    title={#1},
    fonttitle=\sffamily\bfseries\small,
    coltitle=black,
    attach boxed title to top left={yshift=-2mm, xshift=2mm},
    boxed title style={colback=gray!10, boxrule=0pt},
    breakable
}

% Redefine quote to use infobox
\renewenvironment{quote}
  {\begin{infobox}[Обратите внимание]}
  {\end{infobox}}

% --- Hyperlinks (Clean & Modern) ---
\usepackage{hyperref}
\hypersetup{
    colorlinks=true,
    linkcolor=GeoskanBlue,
    filecolor=GeoskanOrange,
    urlcolor=GeoskanBlue,
    pdftitle={Geoskan Pioneer Programming},
    pdfauthor={Сергей Корягин},
    pdfsubject={Educational Material},
    pdfpagemode=UseOutlines,
    bookmarksnumbered=true
}

% --- Table of Contents Styling ---
\usepackage{tocloft}
\renewcommand{\cftsecfont}{\sffamily\bfseries\color{GeoskanBlue}}
\renewcommand{\cftsecpagefont}{\sffamily\bfseries\color{GeoskanBlue}}
\setlength{\cftbeforesecskip}{1em}

% Indentation for sections (Level 1) relative to Part (Level 0)
\setlength{\cftsecindent}{2em}
\setlength{\cftsubsecindent}{4em}

% --- Versioning Logic (Python Managed) ---
\newcounter{buildversion}
\IfFileExists{settings/version.tex}{\input{settings/version.tex}}{\setcounter{buildversion}{0}}

\def\majorversion{1}
\def\minorversion{1} % Bumped version for refactoring
\newcommand{\theversion}{\majorversion.\minorversion.\thebuildversion}

% Time formatting
\makeatletter
\newcommand{\currenttime}{%
  \two@digits{\number\numexpr\time/60\relax}:%
  \two@digits{\number\numexpr\time-60*(\time/60)\relax}%
}
\newcommand{\buildtimestamp}{\two@digits{\number\day}.\two@digits{\number\month}.\number\year\ \currenttime}
\makeatother

% --- 3D Plot Logic (Preserved) ---
\def\centerarc[#1](#2)(#3:#4:#5){%
    \draw [#1] ($(#2)+({#5*cos(#3)}, {#5*sin(#3)})$) arc (#3:#4:#5)%
}
% \newcommand{\tdseteulerzxz}{%
%     \renewcommand{\tdplotcalctransformrotmain}{%
%         \tdplotsinandcos{\sinalpha}{\cosalpha}{\tdplotalpha}
%         \tdplotsinandcos{\sinbeta}{\cosbeta}{\tdplotbeta}
%         \tdplotsinandcos{\singamma}{\cosgamma}{\tdplotgamma}
%         \tdplotmult{\sasb}{\sinalpha}{\sinbeta}
%         \tdplotmult{\sacb}{\sinalpha}{\cosbeta}
%         \tdplotmult{\sacbsg}{\sacb}{\singamma}
%         \tdplotmult{\sacbcg}{\sacb}{\cosgamma}
%         \tdplotmult{\sasg}{\sinalpha}{\singamma}
%         \tdplotmult{\sacg}{\sinalpha}{\cosgamma}
%         \tdplotmult{\sbsg}{\sinbeta}{\singamma}
%         \tdplotmult{\sbcg}{\sinbeta}{\cosgamma}
%         \tdplotmult{\casb}{\cosalpha}{\sinbeta}
%         \tdplotmult{\cacb}{\cosalpha}{\cosbeta}
%         \tdplotmult{\cacbsg}{\cacb}{\singamma}
%         \tdplotmult{\cacbcg}{\cacb}{\cosgamma}
%         \tdplotmult{\casg}{\cosalpha}{\singamma}
%         \tdplotmult{\cacg}{\cosalpha}{\cosgamma}
%         \pgfmathsetmacro{\raaeul}{+\cacg - \sacbsg}
%         \pgfmathsetmacro{\rabeul}{-\casg - \sacbcg}
%         \pgfmathsetmacro{\raceul}{+\sasb}
%         \pgfmathsetmacro{\rbaeul}{+\sacg + \cacbsg}
%         \pgfmathsetmacro{\rbbeul}{-\sasg + \cacbcg}
%         \pgfmathsetmacro{\rbceul}{-\casb}
%         \pgfmathsetmacro{\rcaeul}{+\sbsg}
%         \pgfmathsetmacro{\rcbeul}{+\sbcg}
%         \pgfmathsetmacro{\rcceul}{+\cosbeta}
%     }
% }

% --- Title Page Command ---
\newcommand{\makePioneerTitle}[5]{
    \begin{titlepage}
        \centering
        \vspace*{0.2cm}
        
        % Title & Subtitle (First)
        {\sffamily\Huge\bfseries\color{GeoskanDarkBlue} #1}\\[0.3cm]
        {\sffamily\Large\color{GeoskanBlue} #2}\\[0.5cm]
        
        % Cover Image (Second)
        \begin{tcolorbox}[
            enhanced,
            colback=white,
            colframe=GeoskanBlue,
            boxrule=1.5pt,
            arc=3mm,
            width=0.85\textwidth, % Match the resized image width + padding
            halign=center,
            valign=center,
            boxsep=0pt,
            left=0pt, right=0pt, top=0pt, bottom=0pt,
            shadow={3mm}{-3mm}{0mm}{black!20}
        ]
            \centering
            \resizebox{0.8\textwidth}{!}{ % Reduced from 0.9 to 0.8 (approx 10% reduction)
                \input{images/tikz_drone/drone_cover_code.tex}
            }
        \end{tcolorbox}
        
        \vspace{0.5cm}
        
        % Author Box (Third)
        \begin{tcolorbox}[
            enhanced,
            colback=GeoskanLightGray,
            colframe=GeoskanDarkBlue,
            width=0.85\textwidth,
            arc=0mm,
            boxrule=0pt,
            leftrule=4pt,
            shadow={2mm}{-2mm}{0mm}{black!20}
        ]
            \centering
            \sffamily
            \textbf{Автор:} #3\\
            \textit{#4}
        \end{tcolorbox}
        
        \vspace{0.5cm}
        
        % Footer (Fourth)
        {\sffamily\large #5}
        
        \vspace{0.3cm}
        {\sffamily\footnotesize \textcolor{gray}{Версия документа: \theversion\ (Дата сборки: \buildtimestamp)}}
    \end{titlepage}
}
